\section{熱力学総合演習}

\subsection{熱機関}
いくつかの状態変化を組み合わせ，最初の状態に戻るような過程を{\bf 熱機関}，または{\bf 熱サイクル}と呼ぶ．
最初の状態に戻るので，トータルで温度変化はないことから内部エネルギーの変化量は0[J]である．
熱機関の過程には熱を吸収する吸熱過程と熱を放出する放熱過程がある．
吸熱過程における吸熱量の合計を$Q_{in}$[J]．放熱過程における放熱量を$Q_{out}$とすれば，熱力学第一法則より，熱機関1周で気体が外にする仕事の合計$W_{サイクル}$[J]は
\[W_{サイクル}=Q_{in}-Q_{out}\]
と表せる．

熱機関での1周のエネルギーの変化を，$Q_{in}$[J]の熱を与え，そのうち$W_{サイクル}$[J]分だけ仕事をし，あまりのエネルギー$Q_{out}$[J]を熱として捨てていると解釈できる．
このときの加えた熱量に対してどれだけの仕事をするかを{\bf 熱効率}と呼び，熱効率$e$は
\[\mbox{\boldmath $e=\dfrac{W_{サイクル}}{Q_{in}} =1-\dfrac{Q_{out}}{Q_{in}}$}\]
となる．

また，各過程での$P-V$図における符号付き面積を考えることで，仕事の合計$W_{サイクル}$[J]は熱機関の各過程で囲まれた部分の符号付き面積として表すことができる．
 \begin{figure}[h]
\begin{center}
\includegraphics[width=100mm]{text_p84_fig1.png}
\end{center}
\end{figure}

このとき，時計回りに回る熱機関ならば$W_{サイクル}>0$となり面積と等しいが，反時計回りに回る場合は$W_{サイクル}<0$となり，面積に$-1$をかけた値となる．


\newpage

\begin{figure}[h]
\begin{center}
\includegraphics[width=170mm]{prob_p46.png}
\end{center}
\end{figure}

\newpage

\begin{figure}[h]
\begin{center}
\includegraphics[width=170mm]{prob_p47.png}
\end{center}
\end{figure}

\newpage

\begin{figure}[h]
\begin{center}
\includegraphics[width=170mm]{prob_p136_2.png}
\end{center}
\end{figure}

\newpage

\begin{figure}[h]
\begin{center}
\includegraphics[width=170mm]{prob_p137.png}
\end{center}
\end{figure}


\newpage

\begin{figure}[h]
\begin{center}
\includegraphics[width=170mm]{prob_p138.png}
\end{center}
\end{figure}

\newpage

\begin{figure}[h]
\begin{center}
\includegraphics[width=170mm]{prob_p139.png}
\end{center}
\end{figure}

\newpage

\begin{figure}[h]
\begin{center}
\includegraphics[width=170mm]{prob_p140.png}
\end{center}
\end{figure}

\newpage

\begin{figure}[h]
\begin{center}
\includegraphics[width=170mm]{prob_p141.png}
\end{center}
\end{figure}

\newpage

\begin{figure}[h]
\begin{center}
\includegraphics[width=170mm]{prob_p142_1.png}
\end{center}
\end{figure}
